% Paper 3, §12 — Process supervision after R1: do PRMs still help?
% Anchors: kb/notes/reasoning/process-supervision.md §3-5
%          kb/excerpts/{lightman2023-prm800k,thinkprm2025,deepseek-r1}.md
%          (zheng2025-prm-survey, she2025-r-prm have only stub or
%           abstract-level coverage in the KB as of 2026-05;
%           load-bearing claims are flagged \deepencite.)

\section{Process supervision after R1: do PRMs still help?}
\label{sec:prm-survival}
\label{sec:prm-after-r1}

\Cref{sec:process-reward-models} introduced the verifier leg of the
reasoning triangle on a 2023 generator: \glsterm{process-reward-model}{PRM}-best-of-$N$ dominates
\glsterm{outcome-reward-model}{ORM}-best-of-$N$ at matched $N$ on \glsterm{prm800k}{PRM800K}-grade verifiers, with
min-aggregation as the strongest aggregator and \glsterm{process-reward-model}{PRM}-weighted
\glsterm{self-consistency-2}{self-consistency} as a budget-friendly compromise
\citep[\S 4]{lightman2023-prm800k}. Two years later, the strongest
open-weight \glsterm{reasoning-model}{reasoning model} --- DeepSeek-R1 --- reached its headline
AIME 2024 jump using \emph{terminal verifier rewards only} and
explicitly abstained from any neural \glsterm{reward-model}{reward model}, whether outcome- or
process-based, in the RL loop \citep[\S 2.2]{deepseek-r1}.\footnote{KB
excerpt: \texttt{kb/excerpts/deepseek-r1.md\#sec-2-2-reward}.} The
question this section confronts is the one \cref{sec:process-reward-models}
deferred: in the post-R1 era, do PRMs still earn their training
and inference cost?

The 2025--2026 picture is mixed. PRMs continue to add \glsterm{value}{value} at
\emph{inference} time in a number of measured settings (best-of-$N$
reranking, search guidance) on reasoning-trained policies; their \glsterm{value}{value}
as a \emph{training} signal beyond verifier-only \glsterm{rlvr-3}{RLVR} is contested,
with both diminishing-returns evidence and continued-gain claims
surveyed by \citet{zheng2025-prm-survey}.\deepencite{zheng2025-prm-survey
\S2--3 (excerpt skeleton-only as of 2026-05; PDF acquisition pending)}
This section walks the evidence on both sides, formalises the
generative-\glsterm{process-reward-model}{PRM} dominance claim and its 2026 caveats, and frames the
contradiction that \cref{sec:contradictions} catalogues.

\subsection{R1 is the data point}
\label{sec:prm-survival-r1}

\Citet[\S 2.1, \S 2.2]{deepseek-r1} train \glsterm{r1-zero}{R1-Zero} with the \glsterm{grpo-2}{GRPO}
objective of \cref{eq:grpo} under a rule-based reward:
\begin{equation}
  r_i \;=\; \mathbb{1}\!\left[\, \mathrm{verify}(x, o_i) = \mathsf{true}\,\right]
  \;+\; \alpha\,\mathbb{1}\!\left[\, \mathrm{format}(o_i) = \mathsf{ok}\,\right],
  \label{eq:r1-rule-reward}
\end{equation}
where $\mathrm{verify}(\cdot)$ is an answer-match check on math
problems and a unit-test / compiler check on code, and
$\mathrm{format}(\cdot)$ scores adherence to the
\verb|<think>...</think><answer>...</answer>| schema. No \glsterm{process-reward-model}{PRM} appears
in \cref{eq:r1-rule-reward}; no neural \glsterm{reward-model}{reward model} of any kind appears
in the loop. The DeepSeek authors are explicit about the choice%
\footnote{KB excerpt: \texttt{kb/excerpts/deepseek-r1.md\#sec-2-2-reward}.}
\citep[\S 2.2]{deepseek-r1}:
\begin{quote}
\itshape
[W]e abstain from applying neural reward models — whether outcome-based
or process-based — to reasoning tasks. This decision is predicated on
our observation that neural reward models are susceptible to reward
hacking during large-scale reinforcement learning.
\end{quote}
The headline empirics are well-known: AIME 2024 pass@1 from
$15.6\%$ to $77.9\%$ during RL, and to $86.7\%$ with \glsterm{self-consistency-2}{self-consistency}
at inference time \citep[\S 2.3]{deepseek-r1}.\footnote{KB excerpt:
\texttt{kb/excerpts/deepseek-r1.md\#sec-2-3-aime}.} A \glsterm{reasoning-model}{reasoning model}
that does not place at the AIME human-competitor level cannot be
explained by ``the missing \glsterm{process-reward-model}{PRM}''; the gap is
already closed without one. This is the load-bearing data point that
reframes the post-2023 \glsterm{process-reward-model}{PRM} literature from ``does process supervision
beat outcome supervision?'' (settled by \citet{lightman2023-prm800k} on
the 2023 generator) to ``is process supervision still net-positive on
top of verifier-only \glsterm{rlvr-3}{RLVR}?''

\begin{intuition}
\glsterm{rlvr-3}{RLVR} can be read as \glsterm{outcome-reward-model}{ORM}-with-a-perfect-verifier. The R1 reward
\cref{eq:r1-rule-reward} is precisely the outcome reward
$R^O_\phi(x, z, y) \in \{0, 1\}$ of \cref{eq:orm-loss} with the learned
classifier replaced by an exact symbolic check. Returning to math:
where \citet{lightman2023-prm800k} compared learned-\glsterm{outcome-reward-model}{ORM}-best-of-$N$
against learned-\glsterm{process-reward-model}{PRM}-best-of-$N$, R1 substitutes the learned \glsterm{outcome-reward-model}{ORM} with a
zero-error verifier and pushes the resulting binary signal through \glsterm{grpo-2}{GRPO}
rather than through reranking. The \glsterm{outcome-reward-model}{ORM} half of the 2023 comparison has
been replaced by something strictly stronger; whether the \glsterm{process-reward-model}{PRM} half
still wins requires re-running the experiment on this new generator
class.
\end{intuition}

\subsection{Inference-time PRMs on reasoning-trained policies}
\label{sec:prm-survival-inference}

The cleaner part of the 2025--2026 picture: PRMs as inference-time
re-rankers and search guides on reasoning-trained policies. Two
consistent findings appear across the surveyed literature.

\paragraph{Generative PRMs dominate discriminative ones on
verifiable-step domains.} \Citet{thinkprm2025} report that \glsterm{generative-prm}{ThinkPRM}, a
generative verifier (\cref{eq:thinkprm-gen}) fine-tuned on $1\%$ of
\glsterm{prm800k}{PRM800K} labels and self-generated synthetic verification traces,
beats discriminative PRMs trained on the full \glsterm{prm800k}{PRM800K} on \glsterm{processbench}{ProcessBench},
MATH-500, and AIME 2024 under best-of-$N$ selection and reward-guided
search; out-of-domain it beats the full-\glsterm{prm800k}{PRM800K} discriminative baseline
by $8\%$ on a GPQA-Diamond subset and $4.5\%$ on
LiveCodeBench, and beats \glsterm{llm}{LLM}-as-a-Judge by $7.2\%$ on \glsterm{processbench}{ProcessBench} at
matched verification-token budget \citep[abstract]{thinkprm2025}.%
\footnote{KB excerpt: \texttt{kb/excerpts/thinkprm2025.md\#headline}.}
The mechanism is the one foreshadowed in \cref{sec:prm-thinkprm}: the
verifier's \glsterm{chain-of-thought}{CoT} can apply algebraic-verification, type-checking, and
dimensional-analysis steps that a feed-forward classifier head cannot
internalise from labelled examples alone \citep[\S 5]{thinkprm2025}.%
\footnote{KB excerpt: \texttt{kb/excerpts/thinkprm2025.md\#sec-5}.} The
2025 \glsterm{process-reward-model}{PRM} survey \citep{zheng2025-prm-survey} echoes the framing while
adding a domain caveat: generative PRMs win where each step has a
verifiable certificate (algebraic equality, equation-of-motion balance,
type compatibility) the verifier's \glsterm{chain-of-thought}{CoT} can actually
apply.\deepencite{zheng2025-prm-survey \S3 (excerpt skeleton-only as of
2026-05)}

\paragraph{Reasoning-aware discriminative PRMs close part of the gap at
lower inference cost.} \Citet{she2025-r-prm} train a \glsterm{discriminative-prm}{discriminative PRM}
that emits a short reasoning preamble before scoring; the architecture
remains a verdict head but the input is enriched by a brief
verifier-reasoning prefix. The reported gain is
distribution-shift-robustness rather than peak-benchmark
performance.\deepencite{she2025-r-prm \S3 (excerpt skeleton-only as of
2026-05)} Operationally, this collapses the
\cref{eq:thinkprm-gen}-vs-\cref{eq:prm-signature} dichotomy of
\cref{sec:prm-thinkprm}: as the preamble length $L_{\text{preamble}}$
grows, R-\glsterm{process-reward-model}{PRM} converges to \glsterm{generative-prm}{ThinkPRM}; as it shrinks to zero, R-\glsterm{process-reward-model}{PRM}
recovers the discriminative classifier. The cost
profile for best-of-$N$ scales as $O(N \times L_{\text{preamble}})$,
intermediate between the $O(N)$ classifier and the
$O(N \times L_{\text{verify}})$ generative verifier. \Citet{cao2025-edu-prm}
add an orthogonal axis (entropy-driven step boundaries) which we treat
in \cref{sec:prm-r-edu}; both threads belong to the same finding ---
the post-2023 \glsterm{process-reward-model}{PRM} landscape is more multi-axial than the
\glsterm{prm800k}{PRM800K}-grade-classifier baseline implied.

The \emph{combined} reading: at inference time, the strongest 2026
PRMs are generative, are trained with two orders of magnitude fewer
human labels than \glsterm{prm800k}{PRM800K} (offset by self-generated synthetic
verification traces), and continue to outperform vanilla
discriminative PRMs on competition math benchmarks. Inference-time
process supervision did not vanish post-R1; it shifted form.

\subsection{Training-time PRMs after R1: where the contradiction lives}
\label{sec:prm-survival-training}

The training-time picture is the contested one. Three positions are
attested in 2025--2026.

\paragraph{Position A: PRM-shaped GRPO is theoretically attractive but
empirically risky.} The middle ground of
\cref{sec:prm-r-edu} (and
\cref{eq:r1-rule-reward}'s sparse-binary critique) is the
\emph{denser} reward of substituting or augmenting $r_i$ with a
per-step or per-trace \glsterm{process-reward-model}{PRM} score. Several R1-replication
efforts try this approach
\citep[\S 4]{thinkprm2025}.\footnote{KB note:
\texttt{kb/notes/reasoning/process-supervision.md\#4.2}.} The known
failure mode is \emph{process hacking}: a misaligned \glsterm{process-reward-model}{PRM} rewards
locally-plausible-looking steps without the trace converging on the
correct answer --- the step-level analogue of the outcome-hacking
failure mode that \citet[\S 2.2]{deepseek-r1} cite as their reason for
abstaining from neural reward models entirely. When the \glsterm{process-reward-model}{PRM} aligns
well, dense signal beats sparse signal for credit assignment; when it
misaligns, the policy exploits the gap.

\paragraph{Position B: PRM-best-of-$N$ over R1-class outputs still
adds gains, but smaller ones.} \Citet{thinkprm2025} report continued
\glsterm{process-reward-model}{PRM}-best-of-$N$ gains over reasoning-trained policies on AIME 2024
and MATH-500; \citet{zheng2025-prm-survey} find the same pattern
across surveyed measurements ``until the base model becomes a strong
\glsterm{reasoning-model}{reasoning model} itself,'' at which point the marginal external-\glsterm{process-reward-model}{PRM}
gain shrinks.\deepencite{zheng2025-prm-survey \S4 discussion of
post-R1 measurements (excerpt skeleton-only as of 2026-05)} The
mechanism: as the policy's own internal verification --- the
self-reflection and ``\glsterm{aha-moment}{aha moment}'' behaviours of
\citet[\S 2.3]{deepseek-r1}\footnote{KB excerpt:
\texttt{kb/excerpts/deepseek-r1.md\#sec-2-3-aha}.} --- approaches the
\glsterm{process-reward-model}{PRM}'s signal, the verifier-fidelity gap that the \glsterm{process-reward-model}{PRM} exploits closes.
This does not zero out, but it does erode: the same external \glsterm{process-reward-model}{PRM} that
delivered $20$-point pass@$N$ gains over a 2023 GPT-4 generator on
\glsterm{prm800k}{PRM800K}-MATH delivers smaller gains over an R1-distill at matched $N$.

\paragraph{Position C: process-vs-outcome RL is the wrong question;
the question is which signal is cheap to verify.} The third position,
implicit in \cref{eq:r1-rule-reward}, is that the binary outcome reward
of \glsterm{rlvr-3}{RLVR} is cheap because the verifier is exact (an answer-match check,
a compiler), not because outcome supervision is intrinsically superior
to process supervision. Where exact outcome verification is available
--- competition math, code, formal logic --- the choice between
sparse-but-zero-error and dense-but-noisy is unbalanced; the binary
verifier wins on robustness alone, even if dense signal would
theoretically reduce credit-assignment noise. Where exact outcome
verification is \emph{not} available --- open-domain QA, scientific
reasoning, planning --- neither \glsterm{rlvr-3}{RLVR} nor \glsterm{process-reward-model}{PRM}-shaped \glsterm{rlvr-3}{RLVR} has a clean
empirical answer in 2026. This reframing is consistent with the survey
\citep{zheng2025-prm-survey} but is sharper than its conclusion.

\begin{contradiction}
\textbf{PRMs as a training-time signal beyond verifier-only \glsterm{rlvr-3}{RLVR} ---
unresolved.} The 2023 \glsterm{prm800k}{PRM800K}-era result that \glsterm{process-reward-model}{PRM} > \glsterm{outcome-reward-model}{ORM} at matched
$N$ \citep[Fig.\ 4]{lightman2023-prm800k} was established on a
non-reasoning-trained generator. As the policy class moves to
\glsterm{rlvr-3}{RLVR}-trained reasoning models, the \glsterm{outcome-reward-model}{ORM} in that comparison is replaced
by a zero-error symbolic verifier (\cref{eq:r1-rule-reward}); the \glsterm{process-reward-model}{PRM}
in that comparison must now beat the generator's own internal
verification. \Citet{thinkprm2025} report inference-time gains;
\citet{zheng2025-prm-survey} catalogue mixed empirics, with most
PRMs improving test-time-compute performance ``until the base model
becomes a strong \glsterm{reasoning-model}{reasoning model} itself.'' R1 itself ships with
\emph{no} \glsterm{process-reward-model}{PRM} in the headline RL loop \citep[\S 2.2]{deepseek-r1}.
The \glsterm{process-reward-model}{PRM}-after-R1 question splits into two: (i) at \emph{inference
time}, generative PRMs (\glsterm{generative-prm}{ThinkPRM}, R-\glsterm{process-reward-model}{PRM}) continue to add \glsterm{value}{value}; (ii)
at \emph{training time}, whether \glsterm{process-reward-model}{PRM}-shaped \glsterm{rlvr-3}{RLVR} beats verifier-only
\glsterm{rlvr-3}{RLVR} at matched compute is benchmark-specific and unsettled.
\Cref{sec:contradictions} returns to this as the most consequential
open contradiction of the verifier leg.
\end{contradiction}

\subsection{Process-vs-outcome RL: when does the per-step signal pay?}
\label{sec:prm-survival-cost}

A useful book-keeping question: at fixed total RL compute, when does
the per-step \glsterm{process-reward-model}{PRM} signal pay for the labelling and verifier-inference
cost? Three terms are in tension.

\paragraph{Labelling cost.} Human \glsterm{prm800k}{PRM800K}-style labelling is
$O(\text{problems} \times T \times c_{\text{annot}})$, where
$c_{\text{annot}}$ is annotation cost per step. Math-Shepherd-style
automatic labelling (\cref{sec:prm-mathshepherd}) collapses this to
$O(\text{problems} \times T \times K \times L_{\text{rollout}})$
inference tokens; \glsterm{generative-prm}{ThinkPRM}-style 1\%-\glsterm{prm800k}{PRM800K} + synthetic-verification
collapses it further to $O(\text{problems}_{1\%} \times T \times c_{\text{annot}}
+ \text{synthetic verification traces})$.

\paragraph{Verifier-inference cost in the RL loop.} Per RL step, a
\glsterm{process-reward-model}{PRM}-shaped reward requires $G \times T$ verifier passes (one per step
of each of $G$ rollouts) at cost
\begin{equation}
  C_{\text{PRM-RL}} \;=\; G \times \E[T] \times C_{\text{verify-step}},
  \label{eq:prm-rl-cost}
\end{equation}
versus the verifier-only \glsterm{rlvr-3}{RLVR} cost of $G \times C_{\text{verify-final}}$,
which for an answer-match check is essentially $C_{\text{regex}}
\approx 0$. With a generative verifier, $C_{\text{verify-step}}
\propto L_{\text{verify}}$ tokens; on a long-\glsterm{chain-of-thought}{CoT} trace with $T \in
[10, 100]$ and $L_{\text{verify}} \in [10^2, 10^3]$ tokens, the \glsterm{process-reward-model}{PRM}-RL
verifier-inference cost can dwarf the policy-rollout cost itself.

\paragraph{Signal density.} The benefit side. \glsterm{process-reward-model}{PRM}-shaped \glsterm{grpo-2}{GRPO} replaces
a single binary $r_i$ with a $T$-vector of step scores; under a
well-aligned \glsterm{process-reward-model}{PRM}, the gradient identifies the worst step locally
rather than back-propagating one outcome label across an entire
$10^3$--$10^4$-token trace, which is the credit-assignment story of
\cref{sec:prm-commitment}.

The pay-off condition is roughly that the credit-assignment gain
exceeds the $G \times T \times L_{\text{verify}}$ verifier-inference
overhead and the process-hacking risk. As of 2026-Q2 there is no
published empirical curve that establishes the crossover point on a
reasoning-trained generator at frontier scale; \cref{sec:contradictions}
records this as an open frontier question.

\begin{intuition}
Process supervision is dense supervision, and dense supervision pays
off when the labelling signal is reliable. The 2023 \glsterm{prm800k}{PRM800K}
demonstration was on human labels (very reliable, very expensive);
2024 Math-Shepherd was on rollout-completion-rate labels (cheap,
moderately reliable); 2025 \glsterm{generative-prm}{ThinkPRM} is on a generative verifier's
self-graded signal (cheap, reliability tied to the verifier's own
quality). Returning to math: the \glsterm{process-reward-model}{PRM}-shaped \glsterm{grpo-2}{GRPO} objective replaces
$\hat A_{i,t} = (R_i - \mu)/\sigma$ from \cref{eq:grpo-advantage} with a
per-step advantage $\hat A_{i,t} = R^P_\phi(x, s^{(i)}_{1:t}) - \bar R^P$
that pays off in proportion to $R^P_\phi$'s alignment with ground truth.
When alignment is high, dense beats sparse; when alignment slips, the
policy hacks the gap.
\end{intuition}

\subsection{What the post-R1 evidence actually says}
\label{sec:prm-survival-evidence}

A summary, scoped to what \cref{sec:process-reward-models}'s anchors
plus the 2025--2026 \glsterm{process-reward-model}{PRM} survey support:

\begin{enumerate}
\item \textbf{R1 reached AIME 2024 pass@1 of $77.9\%$ with no \glsterm{process-reward-model}{PRM} in
the RL loop} \citep[\S 2.3]{deepseek-r1}. This single data point
establishes that strong reasoning is achievable from verifier-only
\glsterm{rlvr-3}{RLVR}; the \glsterm{process-reward-model}{PRM} is not a precondition.
\item \textbf{Inference-time generative PRMs continue to outperform
discriminative PRMs} on verifiable-certificate benchmarks
\citep[abstract]{thinkprm2025}. The dominance shifts from a
training-data-volume axis (\glsterm{prm800k}{PRM800K}) to a verifier-compute axis
($L_{\text{verify}}$).
\item \textbf{The marginal \glsterm{value}{value} of an external \glsterm{process-reward-model}{PRM} at inference time
shrinks as the base model becomes a strong \glsterm{reasoning-model}{reasoning model} itself}
\citep{zheng2025-prm-survey}. Mechanism: the policy's internal
verification approaches the \glsterm{process-reward-model}{PRM}'s signal, closing the
verifier-fidelity gap.\deepencite{zheng2025-prm-survey \S4 discussion}
\item \textbf{\glsterm{process-reward-model}{PRM}-shaped \glsterm{rlvr-3}{RLVR} has no clean empirical advantage over
verifier-only \glsterm{rlvr-3}{RLVR} at frontier scale as of 2026-Q2.} The R1 authors
abstain from neural reward models on reward-hacking grounds; the
2025 survey reports mixed empirics.
\item \textbf{R-\glsterm{process-reward-model}{PRM} and EDU-\glsterm{process-reward-model}{PRM}} \citep{she2025-r-prm,cao2025-edu-prm}
demonstrate that the discriminative-vs-generative split is not the
only \glsterm{process-reward-model}{PRM} design axis: reasoning preambles and entropy-driven step
boundaries each yield gains at lower verifier cost than full
generative \glsterm{generative-prm}{ThinkPRM}, suggesting the design space is more
multi-axial than the 2023--2024 literature
implied.\deepencite{she2025-r-prm \S3 and cao2025-edu-prm \S3
(excerpt skeleton-only as of 2026-05)}
\end{enumerate}

What the evidence does \emph{not} settle: whether the verifier leg of
the reasoning triangle is best instantiated as a search-time \glsterm{process-reward-model}{PRM}, a
training-time \glsterm{process-reward-model}{PRM}-shaped reward, or a verifier-only \glsterm{rlvr-3}{RLVR} signal
combined with \glsterm{self-consistency-2}{self-consistency}. The three options coexist in the 2026
frontier model line-up (\cref{sec:frontier-snapshot}) without a
consensus winner.

\subsection{Forward link}
\label{sec:prm-survival-forward}

The \glsterm{process-reward-model}{PRM}-after-R1 question splits cleanly into an inference-time half
(answered: generative PRMs continue to add \glsterm{value}{value}, with diminishing
returns as the policy's own verification improves) and a training-time
half (open: \glsterm{process-reward-model}{PRM}-shaped \glsterm{rlvr-3}{RLVR} vs verifier-only \glsterm{rlvr-3}{RLVR} has no settled
empirical winner). \Cref{sec:frontier-snapshot} provides the
2026 snapshot table of which frontier reasoning models actually ship
PRMs --- and the headline answer is ``almost none in the headline RL
loop, several at inference time.'' \Cref{sec:contradictions} then
collects this section's contradiction alongside the \glsterm{chain-of-thought}{CoT}-faithfulness
question (\cref{sec:cot-faithfulness}), the
\glsterm{rlvr-3}{RLVR}-generalisation-vs-memorisation question
(\cref{sec:rlvr,sec:r1-family}), and the search-vs-RL fixed-compute
tradeoff (\cref{sec:search-vs-rl}) as the open frontier of the
verifier leg.
